\section{From an Arbitrary Code to 172 Mothers}
\label{sec:reduction}
\begin{theorem}\label{thm:main}
There is no binary linear $[43,10,\ge17]$ code.
\end{theorem}
The proof uses E1 and E2 together with the finite computation specified in Section~\ref{sec:computation}. We give the full structural argument here; Appendix~\ref{app:algebra} supplies the auxiliary identities and coordinate details.

\subsection{A forced dual word of weight four}
Suppose that $C_0$ is a $[43,10,\ge17]$ code. Its parity extension $C$ is an even $[44,10,\ge18]$ code: an odd weight at least 17 increases by one, and an even weight at least 17 is already at least 18. Conversely, puncturing any coordinate of an even $[44,10,\ge18]$ code gives a $[43,10,\ge17]$ code. Thus it suffices to exclude $C$.

\begin{lemma}\label{lem:three}
Every three columns of a generator matrix for $C$ are linearly independent.
\end{lemma}
\begin{proof}
If three columns had rank at most two, shortening on them would give a length-41 code of dimension at least eight and distance at least 18. Taking an eight-dimensional subcode contradicts E1, since the Griesmer bound excludes distance at least 19.
\end{proof}
Let $A_w$ and $B_j$ count weight-$w$ words of $C$ and weight-$j$ words of $C^\perp$, respectively. Lemma~\ref{lem:three} implies $B_1=B_2=B_3=0$. Define the length-44 binary Krawtchouk polynomials by
\[
 K_j(w)=\sum_{t=0}^j(-1)^t\binom wt\binom{44-w}{j-t}.
\]
The MacWilliams identities give $\sum_w A_w K_j(w)=2^{10}B_j$. Consider
\begin{align}
 P(w)&=(w-18)(w-20)(w-24)(w-26)\notag\\
 &=\tfrac32K_4(w)+22K_2(w)+\tfrac{403}{2}K_0(w).
 \label{eq:polynomial}
\end{align}
Every allowed nonzero weight is even and at least 18, so $P(w)\ge0$ on these weights. Since $P(0)=224640$ and $B_0=1$,
\[
 224640\le\sum_w A_wP(w)=1536B_4+206336.
\]
Therefore $B_4\ge143/12$, hence $B_4\ge12$. In particular, $C^\perp$ contains a weight-four word.

\subsection{Shortening gives a full-support mother}
\begin{lemma}\label{lem:mother}
Shortening $C$ on the support of a weight-four dual word produces an even full-support $[40,7,18]$ code $K$.
\end{lemma}
\begin{proof}
Permute the four support coordinates to the end. Their generator columns $v_1,v_2,v_3,v_4\in\F^{10}$ satisfy $v_1+v_2+v_3+v_4=0$. Any three are independent, so the projection rank is three. Shortening therefore has dimension $10-3=7$, length 40, and distance at least 18. All its words remain even. If its distance exceeded 18, it would be at least 20, contradicting $g_2(7,20)=42>40$.

Let $U=\langle v_1,v_2,v_3\rangle$. A remaining coordinate would vanish on every shortened word precisely if its generator column belonged to $U$. The seven nonzero vectors of $U$ are the four support columns and their three distinct pair sums. An outside column equal to a support column would give a dependent pair; an outside column equal to a pair sum would give a dependent triple. The zero vector is also excluded by Lemma~\ref{lem:three}. Thus no outside coordinate vanishes, and $K$ has full support.
\end{proof}
By E2 and the verified inequivalence of the 172 actual representatives, $K$ is coordinate-equivalent to one catalogue member. A permutation of its 40 coordinates preserves every distance condition below.

\subsection{Every extension requires compatible cosets}
Choose a basis for the shortened subcode and extend it to a basis of $C$. The image on the final four coordinates is the even-weight subspace of $\F^4$: it has dimension three and is orthogonal to $(1,1,1,1)$. A change of row basis gives
\begin{equation}\label{eq:normal}
 G_C=\begin{pmatrix}
 G_K&0000\\
 a_1&1001\\
 a_2&0101\\
 a_3&0011
 \end{pmatrix}.
\end{equation}
The $a_i\in\F^{40}$ have even weight because $C$ is even. For a vector $a$, set $\dist_K(a)=\min_{c\in K}\wt(a+c)$. Each nonzero combination of one or two of the final three rows has tail weight two; the combination of all three has tail weight four. Consequently the distance condition for \eqref{eq:normal} is exactly
\begin{align}
 \dist_K(a_i)&\ge16 &&(1\le i\le3),\notag\\
 \dist_K(a_i+a_j)&\ge16 &&(i<j),\label{eq:extension}\\
 \dist_K(a_1+a_2+a_3)&\ge14.\notag
\end{align}

Let $E_{40}$ be the even-weight subspace of $\F^{40}$. Since $K\subset E_{40}$, the quotient $E_{40}/K$ has dimension $39-7=32$. Let $F_K$ consist of the cosets satisfying $\dist_K(a)\ge16$. Join distinct cosets $a+K,b+K\in F_K$ when $\dist_K(a+b)\ge16$. Conditions~\eqref{eq:extension} require a triangle in this graph; the three cosets must be distinct because equal cosets have sum-distance zero. The computation finds that all 172 graphs have no edges. This stronger obstruction rules out \eqref{eq:normal} before the triple-sum condition is needed, proving Theorem~\ref{thm:main}.
